% Created 2026-10-01 Thu 19:42
% Intended LaTeX compiler: pdflatex
\documentclass[11pt]{article}

\usepackage[utf8]{inputenc}
\usepackage[T1]{fontenc}
\usepackage{graphicx}
\usepackage{longtable}
\usepackage{wrapfig}
\usepackage{rotating}
\usepackage[normalem]{ulem}
\usepackage{amsmath}
\usepackage{amssymb}
\usepackage{capt-of}
\usepackage{hyperref}
\date{\today}
\title{}
\hypersetup{
 pdfauthor={},
 pdftitle={},
 pdfkeywords={},
 pdfsubject={},
 pdfcreator={},
 pdflang={English}}
\begin{document}

\tableofcontents

\#pragma once
\#include "component.h"
\#include "saveable.h"
\#include <glm/glm.hpp>
\#include <glm/ext.hpp>
\#include <vector>
\#include <string>
\#include "../dflib/dflib.h"

namespace engine \{

class Component;

class Transform : public Saveable \{
public:

Transform();

/**
\begin{itemize}
\item "load" constructor
\item 

\item Called by the level loader to deserialize transforms from json
\item Simply reads the corresponding fields from the json file and populates the class fields
\item 

\item @param value The JSON to deserialize
\end{itemize}
 */
Transform(json value);

/**
\begin{itemize}
\item Same as default constructor, but with the node name set
\item 

\item @param name The desired node name
\end{itemize}
 */
Transform(std::string name);

/**
\begin{itemize}
\item Destroys this transform, all of its components, and then recursively destroys
\item its children
\end{itemize}
 */
\textasciitilde{}Transform();

/**
\begin{itemize}
\item Sets the node's name
\item 

\item @param name The desired node name
\end{itemize}
 */
void SetName(std::string name);

/**
\begin{itemize}
\item Save's the transform's properties to a JSON object
\item 

\item @returns the JSON object representing the transform
\end{itemize}
 */
json Save() override;

/**
\begin{itemize}
\item Load the transform properties saved with `Save()` into the transform
\item 

\item @returns true on success, false on failure
\end{itemize}
 */
bool Load(json value) override;

/**
\begin{itemize}
\item This is the render function called by the editor
\end{itemize}
 */
void ProcessEditorRender();

/**
\begin{itemize}
\item This is the render function called by the game
\end{itemize}
 */
void ProcessRender();

/**
\begin{itemize}
\item Runs once per frame when the game simulation is running
\end{itemize}
 */
void ProcessTick(float dt);

/**
\begin{itemize}
\item This is the initialization function called to initialize
\item components.
\end{itemize}
 */  
void ProcessInit();

/**
\begin{itemize}
\item Add a component to this transform
\item 

\item @param component A pointer to the component to add
\end{itemize}
 \textbf{/  
void AddComponent(Component} component);

/**
\begin{itemize}
\item Get a component on this transform by name
\item 

\item If multiple components have the same name, the first
\item in the array is returned
\item 

\item If no component has the given name, return nullptr
\item 

\item @param name The name of the component to return
\item 

\item @returns A pointer to the component with the given name, or nullptr
\end{itemize}
 \textbf{/  
Component} GetComponent(std::string name);

/**
\begin{itemize}
\item Add a child transform to this transform
\item 

\item If the transform is already a child of this transform
\item log a warning and ignore it
\item 

\item @param tr the transform to add as a child
\item 
\end{itemize}
 */
void AddChild(Transform *tr);

/**
\begin{itemize}
\item Return the child transform that has the given id
\item 

\item If no child has the given ID, return nullptr
\item 

\item @param id The id of the transform to return
\item 

\item @returns A pointer to the transform with the given id, or nullptr
\end{itemize}
 */
Transform *FindTransformByUID(uint64\textsubscript{t} id);

/**
\begin{itemize}
\item Return the attached component that has the given id
\item 

\item If no component has the given ID, return nullptr
\item 

\item @param id The id of the component to return
\item 

\item @returns A pointer to the component with the given id, or nullptr
\end{itemize}
 \textbf{/  
Component} FindComponentByUID(uint64\textsubscript{t} id);

/**
\begin{itemize}
\item Return the child transform with the given name
\item 

\item If more than one child has the given name, return the first in the array
\item 

\item If no child has the given name, return nullptr
\item 

\item @param name The name of the child to return
\item 

\item @returns A pointer to the transform with the given name, or nullptr
\end{itemize}
 */
Transform *FindChildByName(std::string name);

/**
\begin{itemize}
\item Removes the given transform from the children array, and returns it
\item 

\item Does not automatically destroy the transform
\item 

\item If the given transform does not exist in the child array, return nullptr
\item 

\item @param tr The transform to remove
\item 

\item @returns The transform that was removed
\end{itemize}
 \textbf{/
Transform} RemoveChild(Transform *tr);



/**
\begin{itemize}
\item Return the children array of this transform
\item 

\item @returns The children array of this transform
\end{itemize}
 */
std::vector<Transform*> GetChildren();

/**
\begin{itemize}
\item Translates this transform in local coordinates
\item 

\item @param offset The vector to translate the transform by
\end{itemize}
 */
void Translate(glm::vec3 offset);


/**
\begin{itemize}
\item Scales this transform by the given multiplier
\item 

\item @param amount The vector to multiply the scale by
\end{itemize}
 */
void Scale(glm::vec3 amount);

/**
\begin{itemize}
\item Rotates the transform along the given axis by the given amount
\item 

\item @param axis The axis to rotate around
\item @param amount The rotation amount, in degrees
\end{itemize}
 */
void Rotate(glm::vec3 axis, float amount);

/**
\begin{itemize}
\item Returns the forward direction of this transform
\item 

\item @returns the forward direction of the transform
\end{itemize}
 */
glm::vec3 Forward();

/**
\begin{itemize}
\item Returns the backward direction of this transform
\item 

\item @returns the backward direction of the transform
\end{itemize}
 */ 
glm::vec3 Backward();

/**
\begin{itemize}
\item Returns the up direction of this transform
\item 

\item @returns the up direction of the transform
\end{itemize}
 */
glm::vec3 Up();

/**
\begin{itemize}
\item Returns the down direction of this transform
\item 

\item @returns the down direction of the transform
\end{itemize}
 */
glm::vec3 Down();

/**
\begin{itemize}
\item Returns the left direction of this transform
\item 

\item @returns the left direction of the transform
\end{itemize}
 */
glm::vec3 Left();

/**
\begin{itemize}
\item Returns the right direction of this transform
\item 

\item @returns the right direction of the transform
\end{itemize}
 */  
glm::vec3 Right();


/**
\begin{itemize}
\item Calculate the model matrix given the instances position, rotation, and scale
\item 

\item @returns the model matrix corresponding to this instance
\end{itemize}
 */
glm::mat4 GetModelMatrix();

/**
\begin{itemize}
\item Calculates the model matrix with the scale normalized to 1
\item 

\item This is used internally for some editor functionality.
\item 

\item @returns The model matrix corresponding to this instance with a unit scale
\end{itemize}
 */
glm::mat4 GetModelMatrixUnscaled();

/**
\begin{itemize}
\item Calculates the position of this transform in global coordinates
\item 

\item @returns the global position vector of this transform
\item 
\end{itemize}
 */
glm::vec3 GetGlobalPosition();

/**
\begin{itemize}
\item Returns the position of this transform in local coordinates
\item 

\item @returns the local position vector of this transform
\end{itemize}
 */
glm::vec3 GetLocalPosition();

/**
\begin{itemize}
\item Returns the rotation of this transform in global coordinates
\item 

\item @returns the quaternion rotation of this transform in global coordinates
\end{itemize}
 */
glm::quat GetGlobalRotation();

/**
\begin{itemize}
\item Returns the rotation of this transform in local coordinates
\item 

\item @returns the quaternion rotation of this transform in local coordinates
\end{itemize}
 */
glm::quat GetLocalRotation();



/**
\begin{itemize}
\item Returns the scale of this transform in global coordinates
\item 

\item @returns the scale value of this transform in global coordinates
\end{itemize}
 */
glm::vec3 GetGlobalScale();

/**
\begin{itemize}
\item Returns the scale of this transform in local coordinates
\item 

\item @returns the scale value of this transform in local coordinates
\end{itemize}
 */
glm::vec3 GetLocalScale();

/**
\begin{itemize}
\item Returns the name of this transform
\item 

\item @returns the name of this transform
\end{itemize}
 */
std::string GetName();

/**
\begin{itemize}
\item Returns the uid of this transform
\item 

\item @returns the uid of this transform
\end{itemize}
 */
uint64\textsubscript{t} GetUID();


/**
\begin{itemize}
\item Draws the editor UI for this transform
\end{itemize}
 */
void DrawWidget();

glm::vec3 position;
glm::quat rotation;
glm::vec3 scale;
std::string name;
uint64\textsubscript{t} uid;
glm::vec3 rotationEuler; //slaved to the quaternion

Transform *parent = nullptr;
std::vector<Transform *> children;
std::vector<Component*> components;

private:
\};

\}
\end{document}
